\section{Related work}\label{sec:related}

\paragraph{Updating persistence.}
Cohen-Steiner, Edelsbrunner, and Morozov~\cite{CEMS06} introduced vines and vineyards.
They maintain a reduced boundary matrix while two adjacent simplices of the filtration are transposed, in time linear in the number of simplices per transposition.
A continuously changing function is processed as a sequence of such transpositions.
Dey and Hou~\cite{DH21u} carried this approach over to zigzag filtrations.
They identified eight atomic operations, six of which can be executed as efficiently as their non-zigzag counterparts.
Luo and Nelson~\cite{LN21} additionally handle insertions and deletions of cells and batches of changes.
The cost of their scheme scales with the number of elementary changes.
Piekenbrock and Perea~\cite{PP24} replace the long sequences of transpositions that arise when a simplex moves to a new position in the filtration by coarser moves.
Under insertions and elementary collapses, persistence is maintained with annotations~\cite{BCCDW12,DFW14} and by the streaming algorithm for towers of Kerber and Schreiber~\cite{KS19}, and Maria and Oudot~\cite{MO15} compute zigzag persistence online, one insertion or deletion at a time.
In the worst case, every one of these methods spends time polynomial in the size of the complex per elementary change.

Sublinear updates are known in two settings.
For graphs, Dey, Hou, and Parsa~\cite{DHP23} perform a transposition in ordinary graph persistence in $O(\log L)$ time for a filtration of length~$L$, and they update several kinds of bars of graph zigzag persistence in sublinear time.
For time series, the banana trees of Cultrera di Montesano, Edelsbrunner, Henzinger, and Ost~\cite{CEHO24} support insertions, deletions, cuts, and concatenations in time logarithmic in the length of the series plus the number of changes of the augmented persistence diagram.
In both settings the homology that changes is that of a graph, and graph homology reduces to connectivity.
Our results indicate that such bounds do not extend to insertions and deletions in general $2$-complexes unless dynamic rank of sparse matrices admits them too: under the \OMv\ conjecture, amortized update time $n^{1/2-\eps}$ with query time $n^{1-\eps}$ is impossible there, whereas complexes in~$\R^3$ remain easy (Section~\ref{sec:easy-s3}).
A transposition is a different operation from the insertions, deletions, and value changes we consider.
Our lower bounds do not apply to it, and the complexity of vineyards remains open (Section~\ref{sec:conclusion}).

\paragraph{Static complexity of Betti numbers and persistence.}
The standard persistence algorithm~\cite{ELZ02,ZC05} reduces the boundary matrix by Gaussian elimination and takes cubic time in the worst case.
Milosavljevi\'c, Morozov, and \v{S}kraba~\cite{MMS11} showed that persistence, and even zigzag persistence, can be computed in matrix multiplication time.
Chen and Kerber~\cite{CK13} gave an output-sensitive algorithm.
Over finite fields, Wiedemann's method~\cite{Wie86,KS91} computes the rank of an $n\times n$ matrix with $O(n)$ nonzero entries, and hence every Betti number of a complex of fixed dimension, in time $O(n^{2+o(1)})$, as Edelsbrunner and Parsa note~\cite{EP14}; Casacuberta and Kyng~\cite{CK22} also compute a basis of the null space of such a matrix in expected time $O(n^{2.2131})$.
In the other direction, Edelsbrunner and Parsa~\cite{EP14} proved that, over $\Ftwo$ and over prime fields, computing the Betti numbers of a $2$-complex with $N$ simplices has the same complexity as computing the rank of an $N\times N$ $0/1$ matrix with $N$ nonzero entries.
In their construction every column of the matrix becomes a cycle, and every row becomes a sphere with holes, triangulated with $O(h)$ simplices if it has $h$ holes, and glued to the cycles of its nonzero entries.
They also sparsify the matrix to at most three nonzero entries per row and column~\cite[Theorem~4.1]{EP14}, a special case of the sparsification of Yuster~\cite{Yus08}, and they embed their complex in~$\R^4$ by placing its cycles in a plane and every punctured sphere in its own $3$-dimensional half-space from a pencil through that plane, even as a subcomplex of a triangulation of the whole space computed in linear time.
Our gadget follows the same pattern, with the roles of rows and columns exchanged.
What is new is the switches, the decomposition of each punctured sphere into constant-size sockets, which makes every entry update local, and the orientation analysis, which gives the static reduction over every field as well; our embedding of the hard universes in Section~\ref{sec:lower-r4} uses their pencil construction.
This contrasts with complexes in $\R^3$, whose Betti numbers Delfinado and Edelsbrunner~\cite{DE95} compute in near-linear time; Dey and Guha~\cite{DG98} compute homology groups and their generators for arbitrary complexes in $\R^3$.
The contrast between the $3$-sphere and $\R^4$ in Sections~\ref{sec:easy-s3} and~\ref{sec:lower-r4} is the dynamic counterpart of this static picture.
Recent work extends static hardness in other directions: to sufficiently fine additive and to constant-factor multiplicative approximations of Vietoris--Rips and \v{C}ech barcodes~\cite{KM26}, and to the average case, via planted cliques, for estimating Betti numbers and deciding whether homology vanishes~\cite{SSW26}.
Closest to our equivalence on the static side, Ding, Kyng, Probst Gutenberg, and Zhang~\cite{DKPZ22}, building on Kyng and Zhang~\cite{KZ20}, showed that solving linear equations in boundary operators and combinatorial Laplacians of $2$-complexes is as hard as solving general sparse linear equations, by a reduction that preserves sparsity up to polylogarithmic factors.
Our gadget plays the analogous role for rank and Betti numbers; what is new is that the encoding preserves the instance size up to constant factors and is local, so that every change of an entry and every switch of a row or column costs a constant number of simplex updates.
Embeddings in~$\R^3$ also help for Laplacian systems, which can be solved in nearly linear time, parameterized by the first Betti number, for subcomplexes of collapsible complexes embedded in~$\R^3$~\cite{BN22}, whereas for general $2$-complexes they are as hard as general sparse linear equations~\cite{DKPZ22}.
To our knowledge, our results are the first conditional lower bounds for maintaining homology in dimension at least one under updates.

\paragraph{Conditional lower bounds for dynamic problems.}
Henzinger, Krinninger, Nanongkai, and Saranurak~\cite{HKNS15} introduced the \OMv\ conjecture.
They derived from it tight lower bounds for many dynamic problems: amortized bounds for fully dynamic algorithms and worst-case bounds for partially dynamic ones.
Abboud and Vassilevska Williams~\cite{AV14} obtained lower bounds for dynamic problems from other popular conjectures.
\OMv-hardness of dynamic matrix rank over large fields was known.
As van den Brand, Nanongkai, and Saranurak observe~\cite[Theorem~E.5]{vdBNS19full}, for row and column updates it follows from the \OMv\ lower bound for $st$-reachability~\cite{HKNS15} via bipartite matching~\cite{AV14}.
Over every field, including $\Ftwo$, a direct and deterministic reduction from \OuMv\ to dynamic rank follows from the observation that $u^{\top}Mv=1$ exactly when $M[\supp u,\supp v]\ne0$, together with switches or the padding identity~\eqref{eq:padding} (Lemma~\ref{lem:rank-omv}).
Composing it with our equivalence yields the lower bounds for Betti numbers, with a constant number of simplex updates per coordinate of a query vector.
Van~den~Brand et al.\ also proposed hinted variants of \OMv.
These imply worst-case lower bounds for dynamic matrix inverse, determinant, and rank under entry updates that match their algorithms~\cite[Corollary~5.15]{vdBNS19full}, and we discuss what they imply for Betti numbers in Section~\ref{sec:lower-worst}.

\paragraph{Dynamic linear algebra.}
Sankowski maintained the inverse and the determinant of a matrix under entry updates in subquadratic time~\cite{San04}.
Sankowski also reduced dynamic rank to dynamic determinant with randomization~\cite{San07}.
Frandsen and Frandsen~\cite{FF09} studied dynamic matrix rank directly.
Van~den~Brand, Nanongkai, and Saranurak~\cite{vdBNS19} improved the update time for entry updates to $O(s^{1.407})$ in the worst case, for $s\times s$ matrices; for rank the algorithm is randomized.
All these bounds are stated for dense matrices of a given dimension.
Our equivalence identifies dynamic Betti numbers with dynamic rank of sparse matrices, measured by the instance size $N=r+k+m$ of Problem~\ref{prob:rank}.
So an improvement for sparse matrices in this measure would transfer to Betti numbers, and conversely.

\paragraph{Connectivity, low dimension, and bounded treewidth.}
The zeroth Betti number is the number of connected components.
It can be maintained in $O(\log^2 n)$ amortized time by the deterministic algorithm of Holm, de Lichtenberg, and Thorup~\cite{HDT01}, improved by Wulff-Nilsen~\cite{WN13}.
Kapron, King, and Mountjoy~\cite{KKM13} achieve polylogarithmic worst-case time with a Monte Carlo algorithm, and Chuzhoy et al.~\cite{CGLNPS20} achieve deterministic worst-case time $n^{o(1)}$.
The amortized bounds are optimal up to a logarithmic factor: P\u{a}tra\c{s}cu and Demaine~\cite{PD06} proved an $\Omega(\log n)$ cell-probe lower bound per operation for dynamic connectivity, and it transfers to~$\beta_0$, and hence to our algorithms for complexes in~$\R^3$.
Eppstein~\cite{Epp03} maintains a graph embedded on a surface, together with generators of the fundamental group of the surface, in polylogarithmic amortized time per update; this is the closest precedent for our combination of duality and dynamic connectivity.
Delfinado and Edelsbrunner~\cite{DE95} process a filtration of a triangulated $3$-sphere with two union--find structures, one on the complex and one, run backwards, on the dual graph of its complement.
Our algorithm for the $3$-sphere replaces both by fully dynamic connectivity, which is what allows arbitrary insertions and deletions.
Given a tree decomposition of width~$t$ of the bipartite graph of the nonzero entries, Fomin, Lokshtanov, Saurabh, Pilipczuk, and Wrochna~\cite{FLPSW18} compute the rank of an $s\times s$ matrix with $O(t^3s)$ arithmetic operations.
Our dynamic structures store small summaries of groups of columns in a balanced hierarchy (Theorem~\ref{thm:hier}); with a balanced tree decomposition~\cite{BH98}, they show that bounded treewidth of the Hasse diagram suffices for polylogarithmic worst-case updates (Corollary~\ref{cor:tw-betti}).

\paragraph{Output-sensitive and Morse-theoretic methods.}
Several methods aim for a cost governed by the homology rather than by the size of the complex.
Annotations~\cite{BCCDW12,DFW14} maintain a homology basis under elementary inclusions and collapses.
Discrete Morse theory~\cite{For98} and its use for persistence~\cite{MN13} replace a complex by a smaller chain complex on its \emph{critical cells}, the cells left unmatched by an acyclic matching between cells and their facets; their number is at least the total Betti number.
The hard instances of Theorem~\ref{thm:acyclic} are connected and have $\beta_1=\beta_2\le1$, so no method whose cost is governed by the homology of the complex, rather than its size, can escape our lower bounds.
This argument does not cover methods whose cost depends on the number of critical cells of a given matching, which can be much larger than the total Betti number.
